\chapter{Zusammenfassung und Ausblick}
\label{ZusammenfassungAusblick}
% GEÄNDERT: Vorlagentext durch Gerüst ersetzt

\section{Beantwortung der Forschungsfrage}
\label{sec:ff-antwort}
% TODO: Teilfragen TF1-TF4 nacheinander "abhaken" (Leitfaden), Grad der Beantwortung aus dem Profil ableiten

\section{Limitationen}
\label{sec:limitationen}
% TODO: Limitationen des Ansatzes (Einzelsystem, Korpus nur Customizing-Tabellen, Asymmetrie Standard-/Kundenobjekte, Referenzinhalt des Standards);
% ggf. hier knapp: bewusst nicht verfolgte Lösungswege (Meeting-Notiz: Abgrenzung "wenn dann erst später")

\section{Ausblick}
\label{sec:ausblick}

% VERSCHOBEN aus 1.2 (Entscheidung R5) und GEÄNDERT: von der Ankündigung ("wird ... untersucht") in einen Ausblick umformuliert
Eine naheliegende Erweiterung besteht darin, den Ansatz von der Abweichung gegenüber dem Herstellerstandard auf den Vergleich zweier ursprünglich gleichartiger Systeminstanzen zu übertragen, um Abweichungen zwischen einem Ausgangs- und einem abgeleiteten System aufzudecken. % NEU: Folgesatz zur Begriffsschärfung (Textbaustein 10.09.2026)
Ein solcher Vergleich misst allerdings nicht die Abweichung vom Auslieferungszustand, sondern die Veränderung innerhalb einer Systemlandschaft, sodass der in dieser Arbeit verwendete Begriff der Abweichung hierfür zu erweitern wäre.
